\subsection{The full metabelian section in the specified class is cuspidal}

\begin{statement}{A.169}{Proved}
{Complete metabelian classification in the actual degree-eight Kummer class}
Every continuous section of the full geometrically metabelian
extension \(E_U/\Pi''\) which locally matches actual point or cusp
sections and satisfies
\[
                 2s^*\kappa_8(t)+s^*\kappa_8(1-t)=[64a^2]
\tag{A.169.1}
\]
is geometrically conjugate, in that same full metabelian extension,
to the image of an actual rational cusp section of \(U\).

More precisely, choose a global fixed sheet of \(D^\circ\), and
regard the resulting section as a section of \(E_D/\Pi''\).
Modulo conjugacy by \(H/\Pi''\), its class is the image of exactly
one full cusp class on \(D^\circ\), specified by
\[
             P\in\{P_0,P_1,P_{1/2},P_\infty\},\qquad
             \tau\in\widehat{\Q^\times}
\tag{A.169.2}
\]
relative to the four parameters in A.167.  Thus the map from these
locally matching full metabelian classes to their affine abelian
cover classes in A.167 is a bijection.  Downstairs their exact set of
geometric conjugacy classes is
\[
             \coprod_{P\in\{P_0,P_1,P_{1/2},P_\infty\}}
                         \widehat{\Q^\times}/\{\pm1\}.
\tag{A.169.3}
\]
This statement concerns the specified combined class (A.169.1).
It imposes no assertion of classification on other Kummer classes.
\end{statement}

\begin{proof}
\emph{Completed Kummer detection over every number field.}
We first prove the precise local-global fact needed for the boundary
fields of the covers.  For any number field \(k\), localization is
injective on
\[
 H^1_{\rm cts}(k,\widehat\Z(1))
        =\widehat{k^\times}
       \longrightarrow\prod_{v\ {
m finite}}
                                  \widehat{k_v^\times}.
\tag{A.169.4}
\]
The same assertion holds when the product is taken over a set of
finite places of Dirichlet density one.  The continuous Kummer
identification and its local maps follow from the finite
zero-cochain argument of A.167.10, which works over every such field.

Fix a coefficient prime \(\ell\), and represent a class in the
localization kernel by a continuous cocycle
\(c:G_k\to\Z_\ell(1)\).  At level \(\ell^r\), let
\(a_r\in k^\times\) represent its finite Kummer class and put
\[
 F_r=k(\mu_{\ell^r}),\qquad
 L_r=F_r(a_r^{1/\ell^r}),\qquad
 N_r=\operatorname{Gal}(L_r/F_r).
\tag{A.169.5}
\]
A nonidentity element of \(N_r\) fixes no root of
\(X^{\ell^r}-a_r\): it fixes the roots of unity and translates
every root by the same nonidentity multiplier.  Its conjugacy class
in the finite Galois group over \(k\) therefore gives, by
Chebotarev, a positive-density set of finite places at which the
polynomial has no local root.  The comparison is the literal
unramified Frobenius permutation, after omitting the finitely many
places dividing \(\ell a_r\) or ramified in \(L_r\).
This contradicts the assumed local zero class, even with a
density-one set of places.  Hence \(N_r=1\) for every \(r\)
\cite[Theorem 8.31]{AgentMilneANT}.

Put \(k_\infty=k(\mu_{\ell^\infty})\).
The finite restriction of \(c\) is zero over \(F_r\).
On \(G_{k_\infty}\) the action on \(\mu_{\ell^r}\) is trivial,
so a zero cohomology class is the zero cocycle there.  Thus \(c\)
vanishes literally on \(G_{k_\infty}\), at every finite level and
therefore in \(\Z_\ell(1)\).  It descends continuously to the
abelian group
\(\Gamma=\operatorname{Gal}(k_\infty/k)\subset\Z_\ell^\times\).
Choose a Tate basis, so its action is multiplication by the
cyclotomic character \(\chi\).  The cocycle identity for the two
orders of a commuting pair gives
\[
    (\chi(\sigma)-1)c(\gamma)
                  =(\chi(\gamma)-1)c(\sigma).
\tag{A.169.6}
\]
The cyclotomic tower over a number field is infinite.  Choose
\(\sigma\) with \(\chi(\sigma)\ne1\), and set
\(z=c(\sigma)/(\chi(\sigma)-1)\in\Q_\ell(1)\).
Equation (A.169.6) says
\(c(\gamma)=(\chi(\gamma)-1)z\) for every \(\gamma\).
At any one of the finite places where the class vanishes there is
\(z_v\in\Z_\ell(1)\) giving the same coboundary locally.
The local cyclotomic action is nontrivial: at residue characteristic
different from \(\ell\), a Frobenius has value the residue-field
cardinality, which is not one in \(\Z_\ell\); at a place over
\(\ell\), the local cyclotomic tower is infinite.  Thus
\((\chi(\gamma_v)-1)(z-z_v)=0\) for a nonzero scalar, and
\(z=z_v\) is integral.  The cocycle is globally a coboundary.
Applying this argument to each coefficient prime proves (A.169.4).
Continuous cocycles and coboundaries for the product
\(\widehat\Z(1)=\prod_\ell\Z_\ell(1)\) are exactly their
coordinatewise continuous cocycles and coboundaries.

\emph{Restriction and the exact nonsplit boundary module.}
We record the finite-extension descent map used below.  If \(L/k\)
is finite Galois, then
\[
 H^1_{\rm cts}(k,\widehat\Z(1))
       \xrightarrow{\ \operatorname{res}\ }
 H^1_{\rm cts}(L,\widehat\Z(1))^{\operatorname{Gal}(L/k)}
                         \quad\text{is an isomorphism}.
\tag{A.169.7}
\]
The relevant invariants are zero:
\(\widehat\Z(1)^{G_L}=0\), since the roots of unity of the
number field \(L\) form a finite group and the inverse Tate
transitions eventually annihilate every fixed finite level.
Here is the actual continuous cocycle inverse to restriction.
Let \(c\) on \(G_L\) represent an invariant class.  For every
\(g\in G_k\), there is a unique \(b_g\in\widehat\Z(1)\)
satisfying
\[
 g\,c(g^{-1}hg)-c(h)=(h-1)b_g\qquad(h\in G_L).
\tag{A.169.8}
\]
Existence is invariance of the class; uniqueness is the zero
invariant group.  The two ways of applying (A.169.8) give
\(b_{gg'}=b_g+g b_{g'}\), and for \(g\in G_L\) its value
is \(c(g)\).  On each of the finitely many cosets of \(G_L\)
this formula is continuous.  It is therefore the required global
cocycle.  A class restricting to zero is zero by the same argument
after a coboundary adjustment.  This proves (A.169.7).  The same
proof gives restriction injectivity for every module whose
\(G_L\)-invariants vanish.

Let \(C/k\) be one of the proper curves used below, let \(B_C\)
be its finite geometric boundary, and let \(P\in B_C(k)\).
Its actual boundary module is
\[
 M_C=\frac{\widehat\Z(1)[B_C]}{
                 \widehat\Z(1)\sum_{Q\in B_C}e_Q}.
\tag{A.169.9}
\]
The rational point \(P\) gives a \(G_k\)-equivariant
identification
\[
 M_C\simeq\widehat\Z(1)[B_C\setminus\{P\}],
           \qquad e_P=-\sum_{Q\ne P}e_Q.
\tag{A.169.10}
\]
No splitting of the other boundary points over \(k\) is assumed.
Suppose a class \(\alpha\in H^1(k,M_C)\) has, at every finite
place, an image in the actual inertia line
\(H^1(k_v,\widehat\Z(1)e_P)\).  Choose a finite Galois field
\(L/k\) splitting \(B_C\).  In (A.169.10), its restriction has
one completed Kummer coordinate \(\alpha_Q\in\widehat{L^\times}\)
for each \(Q\ne P\).  At every finite place of \(L\) these
coordinates are all equal: locally the line \(e_P\) is exactly
the negative diagonal.  Equation (A.169.4) over \(L\) gives
\(\alpha_Q=\beta\) globally for one common class \(\beta\).
The restricted vector is Galois invariant as a cohomology class.
Since its coordinates are equal and Galois acts by permuting those
coordinates and by its usual field action, \(\beta\) is invariant.
Equation (A.169.7) descends it to a unique
\(\beta_0\in\widehat{k^\times}\).

Now \(\alpha+\kappa(\beta_0)e_P\) restricts to zero over
\(L\).  The invariants \(M_C^{G_L}\) are zero, again because
the split coordinates are Tate modules of roots of unity.
Restriction injectivity, proved after (A.169.8), therefore gives
\[
                  \alpha=-\kappa(\beta_0)e_P.
\tag{A.169.11}
\]
This is the needed inertia-line recovery for the specified permutation
module.  Its proof supplies the actual scalar descent and does not
invoke a local-global principle for arbitrary tori.

\emph{Only the actual abelian covers of the tripod are used.}
Lift the section in the statement to \(E_D/\Pi''\) by its global
fixed sheet.  A.167 supplies the unique \((P,\tau)\), and choose
the corresponding full cusp section \(\sigma:G_\Q\to E_D\).
Lift its affine-abelian conjugator and conjugate the given section,
so their images in \(E_D/\overline{[H,H]}\) are literally equal.
In particular their geometric abelian images in the tripod are
equal.  Write this common completed pair as \((A,B)\).

For \(n=8j!\), put
\[
 H_n=\ker\bigl(\Pi\longrightarrow\Pi^{\rm ab}/n\Pi^{\rm ab}\bigr),
             \qquad H_n'=\overline{[H_n,H_n]}.
\tag{A.169.12}
\]
These are the subgroups of A.68.  They satisfy
\(H_n\subset H\), \(\Pi'\subset H_n\), and
\(\Pi''\subset H_n'\).  In the finite arithmetic quotient
\(E_U/H_n\), take the common section given by \(s\) and
\(\sigma\), and let \(E_n\subset E_U\) be its full preimage.
Then \(E_n\subset E_D\), its geometric group is \(H_n\),
and both sections lift to its indicated quotients.  This defines an
actual finite etale cover \(Y_n\to D^\circ\to U\), with all
arithmetic subgroup maps fixed.  For \(n\mid m\), these
subgroups and covers have their actual inclusion and transition map.
In particular arbitrary abelian covers of \(D^\circ\), whose
commutator quotient need not receive \(\Pi/\Pi''\), are not
substituted for (A.169.12).

The covers and their map to \(D\) can be written literally.
Choose rational representatives \(A_n,B_n\) of the coordinates
of \((A,B)\) modulo \(n\).  With a choice of the fixed sheet,
\[
 \begin{gathered}
 Y_n:\quad A_nu^n+B_nv^n=1,\quad uv\ne0,
                  \qquad t=A_nu^n,\\
 W=c_nu^{n/4}v^{n/8},\qquad
 T=\frac12-\frac{4aW^4}{t},\qquad
 c_n^8=\frac{A_n^2B_n}{64a^2},\quad c_n\in\Q^\times.
 \end{gathered}
\tag{A.169.13}
\]
The last rational root exists by (A.169.1) and \(8\mid n\).
The two choices of its sign correspond to the two rational sheets
in A.167; choose the one giving the selected map to \(D^\circ\).
The Kummer equations make the pullback along \(\sigma\) a
trivial finite torsor, and its fixed sheet is exactly the arithmetic
subgroup just defined.  The displayed monomial has eighth power
\(t^2(1-t)/(64a^2)\), and its geometric character has order eight;
the inverse formula for \(T\) is the one in A.139.
Thus (A.169.13) is the actual finite etale map of degree \(n^2/8\)
to \(D^\circ\), not only a comparison of quotient groups.
For adjacent levels \(n\mid m\), the actual transition fixed by
\(E_m\subset E_n\) is
\((u_m,v_m)\mapsto(ru_m^{m/n},sv_m^{m/n})\), with
\(r^n=A_m/A_n\) and \(s^n=B_m/B_n\).
These constants are rational: each quotient of the two indicated
root functions has a constant \(n\)-th power and the geometrically
connected curve has constant field \(\Q\).
Choosing the compatible coordinate presentation gives
\(c_m=c_nr^{n/4}s^{n/8}\), so these maps commute with
(A.169.13).  Further transitions are their compositions, retaining
the fixed arithmetic sheets supplied by the inclusions.

The proper smooth model \(C_n\) is the plane Fermat twist
\(A_nU^n+B_nV^n=Z^n\).  Its gradient proves geometric
smoothness, and its genus is \((n-1)(n-2)/2\).
It has a rational boundary point supplied by the lifted full cusp
\(\sigma\).  It is also visible on these equations.  At a cusp
over zero one may take \(B_n=1\) and \((U:V:Z)=(0:1:1)\);
over one take \(A_n=1\) and \((1:0:1)\); over infinity take
\(B_n=-A_n\) and \((1:1:0)\).  These choices use precisely
the actual abelian coordinates of a cusp section.  The finite
fixed-sheet lift chooses its particular boundary point if the
coordinate presentation is changed.  Fix such a rational boundary
point as the Abel--Jacobi origin for \(C_n\).

\emph{The proper image on each cover is one rational boundary point.}
The full metabelian section maps to the arithmetic geometrically
abelian quotient for \(Y_n\), since \(\Pi''\subset H_n'\).
Compactifying and applying its Abel--Jacobi map gives a continuous
class
\[
       \eta_n\in H^1_{\rm cts}(\Q,T\operatorname{Jac}(C_n)),
                \qquad (\eta_{n,m})_m\in
                         \widehat{\operatorname{Sel}}(\Q,
                                      \operatorname{Jac}(C_n)).
\tag{A.169.14}
\]
The local Selmer comparisons here are actual finite-cover maps.
A quotient conjugator lifts geometrically; the common finite sheet
then makes the corresponding full local point section lift to a
point of the finite etale fibre, or a cusp section lift to its
rational boundary point and tangent.  The real fixed-sheet or
inertia-orbit argument is the same as in A.167.  After passing to
the proper curve this is an actual local point at every finite
place and a real component at infinity.

Every such proper local point lies in the finite subscheme
\[
                         Z_n=(C_n\longrightarrow\overline D)^{-1}(P).
\tag{A.169.15}
\]
Indeed its image on the selected \(D\)-sheet is the fixed point
\(P\) of A.163.  The whole fibre in (A.169.15) is in the
boundary of \(Y_n\).  Its points need not be rational, and no such
assumption is made.  Apply Stoll's Theorem 3.11 to its Abel--Jacobi
image.  That exact theorem allows an arbitrary finite subscheme of
an abelian variety over a number field: an inverse-Selmer class
whose local values lie there at density one is the class of a
rational point of that subscheme
\cite[Theorem 3.11, p.~11 of the author version]{LegStoll2007}.
All finite places satisfy the hypothesis here.  The Abel--Jacobi
map of the positive-genus curve is injective on geometric points,
so it gives one \(Q_n\in Z_n(\Q)\) with
\(\eta_n=\kappa([Q_n-Q_n^0])\).  Equality holds for the
continuous class by the finite zero-cochain compactness comparison.
Full local Kummer injectivity then makes every selected finite
local centre equal to \(Q_n\).  No Tate--Shafarevich hypothesis
on these larger Jacobians has been introduced.

\emph{Recovering the full affine cusp on every cover.}
The exact compactification sequence for \(Y_n\) is
\[
 0\longrightarrow M_n\longrightarrow H_n^{\rm ab}
           \longrightarrow T\operatorname{Jac}(C_n)\longrightarrow0,
 \qquad
 M_n=\widehat\Z(1)[C_n\setminus Y_n]\big/
                    \widehat\Z(1)\sum e_Q.
\tag{A.169.16}
\]
It is the actual homology sequence: the kernel is generated by the
boundary meridians with their sole sum relation, and Galois acts by
the permutation of the boundary points and the cyclotomic character.
This can be checked after splitting the boundary by the punctured
surface presentation, and it descends with exactly that permutation
action.  Equivalently it is the boundary-torus sequence in A.68.4--5,
with its actual generalized-Jacobian and compactification maps.

Use a rational local parameter at the rational point \(Q_n\) and
its full tangential splitting.  Lift the proper geometric conjugator
in (A.169.14).  The difference from this splitting is then a class
\(\alpha_n\in H^1(\Q,M_n)\).  The map from this cohomology
group to \(H^1(\Q,H_n^{\rm ab})\) is injective, and so is its
local version at every finite place: the compact Jacobian's Tate
invariants vanish, since its torsion over each finite local field
is finite.  This is exactly the cocycle proof of A.167.8--9 and
does not depend on the rank of \(\operatorname{Jac}(C_n)(\Q)\).
The actual local image must be a cusp at \(Q_n\), since an
ordinary \(Y_n\)-point cannot compactify to that boundary point.
Consequently \(\alpha_n\) is locally in the inertia line of
\(Q_n\).  Apply (A.169.9)--(A.169.11), with this rational
distinguished point and all the other boundary orbits retained.
It gives a global completed tangent parameter and an actual full
cusp section \(\sigma_n\) on \(Y_n\) whose image equals the
given section in \(E_n/H_n'\), up to geometric conjugacy.

Push \(\sigma_n\) through the actual finite cover to
\(D^\circ\).  It is a full cusp section at \(P\): on completed
local parameters the map is \(s_P=c s_{Q_n}^{e}u(s_{Q_n})\),
with \(c\in\Q^\times\), \(u(0)=1\), and the actual
ramification index \(e\); its full tangent class is
\(c\tau_n^e\), as proved by formal roots in A.167.  Its image
in \(E_D/\overline{[H,H]}\) equals that of \(s\), hence
equals that of the fixed \(\sigma\).  The selected-cover
bijection of A.167 now forces exactly the same \(P\) and the
same completed tangent parameter \(\tau\).  Full cusp sections
with these data are conjugate in the actual decomposition group
after their geometric paths to \(P\) have been identified: their
two cocycles in \(\widehat\Z(1)\) differ by a coboundary.
Their paths differ by an element of \(H\).  Thus, for every
\(n=8j!\), the section \(s\) is conjugate to the fixed full
\(\sigma\) in \(E_U/H_n'\), by an element of \(H/H_n'\).

\emph{One conjugator in the full metabelian quotient.}
For each such \(n\), let
\[
 \mathcal C_n=\left\{g\in H/\Pi'':
      g s(\gamma)g^{-1}=\sigma(\gamma)
       \text{ in }E_U/H_n'\text{ for every }\gamma\in G_\Q\right\}.
\tag{A.169.17}
\]
This is a nonempty closed subset of the compact group \(H/\Pi''\).
Nonemptiness follows from the preceding conjugator and the surjection
\(H/\Pi''\to H/H_n'\); closedness is the intersection of the
equality conditions in a profinite Hausdorff group.  For \(n\mid m\)
the actual transition gives \(\mathcal C_m\subset\mathcal C_n\).
Compactness therefore supplies one element in their intersection.

The exact kernel identity is
\[
                         \bigcap_{j\ge1}H_{8j!}'=\Pi''.
\tag{A.169.18}
\]
To recall its proved topology from A.68, a continuous finite abelian
character of \(\Pi'\) detecting an element outside \(\Pi''\)
extends to \(\Pi'W\), where \(W\) is open normal and
\(W\cap\Pi'\) lies in its kernel.  Commutators with \(W\)
are killed, so this is a homomorphism.  The open image of
\(\Pi'W\) in \(\Pi^{\rm ab}\) contains \(n\Pi^{\rm ab}\)
for a sufficiently divisible \(n=8j!\).  Thus that character
kills \(H_n'\), and the detected element is outside \(H_n'\).
An element outside \(\Pi'\) is already outside every \(H_n'\).
The reverse inclusion is \(\Pi''\subset H_n'\).
This proves (A.169.18).  The common conjugator in (A.169.17)
therefore gives equality in \(E_U/\Pi''\) itself.

It follows that every selected-cover class is precisely the full
metabelian image of the unique cusp class \((P,\tau)\).
Conversely every such full cusp class is locally matching and
satisfies (A.169.1).  On the selected cover, distinct data cannot
be conjugate in the full metabelian quotient because their affine
images are distinct by A.167.

For the downstairs conjugacy, both selected-cover sections have the
same finite arithmetic image \(E_D/H\) in \(E_U/H\).
A geometric conjugator in \(\Pi/\Pi''\) therefore has an image
\(d\in\Pi/H=C_8\) centralizing that finite section.
The selected-sheet action on the literal \(\mu_8\)-torsor is
cyclotomic, so \(d\in C_8^{G_\Q}=\{0,4\}\).
The zero class is already the selected-cover conjugacy.  The other
class is the actual rational deck automorphism
\((T,W)\mapsto(T,-W)\); it fixes all four boundary points and
negates each specified local parameter.  It takes \(\tau\) to
\(-\tau\), and its action is realized by full geometric path
conjugation, so supplies that identification in the full metabelian
quotient as well.  No other class can occur.  This proves exactly
(A.169.2)--(A.169.3).

The exact Fermat-to-\(D\) maps, their ramification and genus
identities, and the scalar cocycle and boundary-permutation maps are
checked in \path{code/harmonic_full_metabelian_certificate.py}.
The inverse-limit conclusion uses the proved compact conjugator
argument, rather than an inference from uncoordinated section classes.
\end{proof}
